\HMTNarrativeHeading{Action, aire et information}

Le paramètre de Barbero--Immirzi est présenté sur les antécédents
d’action, d’orientation et d’incidence déjà construits. Les coordonnées
angulaires et les canaux conjugués agissent avec l’inclusion
marquée entre configurations hexadiques et octadiques. La fonctionnelle
reçoit ces données et détermine son évaluation. La relation entre
arithmétique et géométrie demeure visible dans chacun de ses arguments.

La structure d'incidence organise quels éléments sont distingués et comment ils se prolongent. Une hexade est un support de six positions du plan combinatoire ; son élévation utilise une dodécade et un alignement déclarés. L'octade résultante appartient à la réalisation binaire. Les lecteurs d'aire conservent la marque qui donne sens à cette correspondance. Les figures aident à suivre le transport des ensembles et de leurs intersections.

L’action et l’horloge permettent ensuite d’exprimer l’énergie associée
à une lecture. La constante de Boltzmann intervient dans la transduction
thermique ; le paramètre de Barbero et l’unité aréale interviennent dans
l’incrément géométrique. La constante gravitationnelle entre dans la
réalisation qui relie l’action et l’aire de Planck. Lors de la composition de ces
expressions, une même section d’action peut s’annuler dans le
quotient entre l’énergie et l’incrément d’aire. L’annulation exhibe
une dépendance structurelle concrète.

La loi de changement d’état conserve aussi le terme d’entropie
relative et la référence par rapport à laquelle la comparaison est réalisée.
L’information d’une observation réduite dépend des corrélations
que maintient l’état conjoint. Cette distinction prolonge l’argument
de l’archive : le contenu de la totalité et celui d’un lecteur partiel
sont reliés au moyen d’opérations explicites.

La théorie de conservation acquiert ainsi une application physique déterminée. L'état, son orientation, l'aire qui lui est attribuée et la lecture thermique forment une composition aux antécédents identifiés. Le développement conserve ces antécédents et leurs preuves. La synthèse ultérieure réunira cette composition avec les bilans opératoriels et avec l'exemple d'intrication calculable.

